\vspace{4pt plus 4pt minus 1pt}\noindent{\kbvsecfont\fontsize{11pt}{14.5pt}\selectfont\color{kbvaccentdark} (೨)\hspace{0.45em}ಪ್ರಾಸ:}\hspace{0.5em}ಯಾವುದೇ ಪದ್ಯದ ಪ್ರತಿಯೊಂದು ಪಾದದ ಮೊದಲನೆಯ ಮತ್ತು ಎರಡನೆಯ ಸ್ವರಗಳ ನಡುವೆ. ಒಂದೇ ಜಾತಿಯ ವ್ಯಂಜನಾಕ್ಷರ ಬರುವುದನ್ನು “ಪ್ರಾಸ” ಎನ್ನುತ್ತೇವೆ. ಮೊದಲನೆಯ ಪಾದದ ಪ್ರಾಸಾಕ್ಷರದ ಹಿಂದಿನ ಅಕ್ಷರ ಹ್ರಸ್ವಸ್ವರವಾಗಿರಲೀ, ದೀರ್ಘಸ್ವರವಾಗಿರಲೀ, ಅನುಸ್ವಾರವಾಗಿರಲೀ ವಿಸರ್ಗವಾಗಿರಲೀ ಆಗಿದ್ದ ಪಕ್ಷದಲ್ಲಿ ಉಳಿದೆಲ್ಲ ಪಾದ ಗಳಲ್ಲಿಯೂ ಅದೇ ಲಕ್ಷಣ ಕಂಡುಬರುವುದು. ಇದು ಪ್ರಾಸ ನಿಯಮ.\par

\vspace{3pt plus 2pt}\noindent\hspace*{1.2em}\kbvlabel{ಉದಾ}\par\vspace{1pt}

\indent ಮಡದಿ ನಿಜ ನಿಳಯವ ಬಿಟ್ಟಡಿ\par

\kbvsechead{ಗಡಿಗೆ ಪರಗೃಹದೊಳಗೆ ಬಾಯನು :}

\begin{kbvverse}
\noindent ಬಡಿದು ಮನೆಮುಕೆನಾರ್ತೆಯೆನ್ನದೆ ಬೀದಿಗಲಹವನು ।
\end{kbvverse}

\indent ಒಡರಿಚುವ ಸತಿಯೊಬ್ಬನುಂಬೆಂ\par

\begin{kbvverse}
\noindent ದೆಡಹಿ ಕಾಣದ ದಿಟ್ಟೆ ಹತ್ತನು ।
\end{kbvverse}

\indent ಹಡೆದೊಡೆಯು ವರ್ಜಿಸುವುದುತ್ತಮ ಪುರುಷರುಗಳೆಂದ\par

\noindent\hfill{\fontsize{9.5pt}{12.5pt}\selectfont\color{kbvgrey}—ಕರ್ಣಾಟ ಭಾರತ ಕಥಾಮಂಜರಿ: ಉದ್ಯೋಗಪರ್ವ}\par\vspace{3pt plus 2pt}

\indent ಮಡ = ಮ್ + ಅ + ಡ್ + ಅ; ಒಡ = ಒ + ಡ್ + ಅ\par

\indent ಗಡಿ = ಗ್ + ಅ + ಡ್ + ಇ; ದೆಡ = ದ್ + ಎ + ಡ್ + ಅ\par

\indent ಬಡಿ = ಬ್ + ಅ + ಡ್ + ಇ; ಹಡೆ = ಹ್ + ಅ + ಡ್ + ಎ\par

\begin{kbvnote}
\noindent{\kbvsecfont ವಿ.ಸೂ. :}\hspace{0.5em}“ರಗಳೆ”ಗಳಲ್ಲಿ ಪಾದಾಂತ್ಯ (ಸಾಲಿನ ಕೊನೆ) ದಲ್ಲೂ ಪ್ರಾಸ ನಿಯಮ ಕಂಡುಬರುವುದು ; ಇದೊಂದು ವಿಶಿಷ್ಟ ಲಕ್ಷಣವೇ ಸರಿ.
\end{kbvnote}

\vspace{4pt plus 4pt minus 1pt}\noindent{\kbvsecfont\fontsize{11pt}{14.5pt}\selectfont\color{kbvaccentdark} (೩)\hspace{0.45em}ಗಣ:}\hspace{0.5em}ಪದ್ಯದ ಪ್ರತಿಪಾದದಲ್ಲಿರುವ ಅಕ್ಷರಗಳ ಗುಂಪಿಗೆ “ಗಣ” ಎಂದು ಹೆಸರು. ಗಣಗಳಲ್ಲಿ ಪ್ರಧಾನವಾಗಿ ಮಾತ್ರಾ ಗಣಗಳು ಮತ್ತು ಅಕ್ಷರ ಗಣಗಳು ಎಂದು ಎರಡು ವಿಧ.\par
